\chapter{UML and OOP}
\label{chap:c++newfeatures}
\vspace{10.5cm}

%%% UML 
\section{UML}
    %% Class Diagram
    \subsection{Class Diagram}
    In Unified Modeling Language, a class diagram \highlightit{describes 
    the static structure of a system} by showing the system's classes, 
    their attributes, operations (methods) and the relationship among 
    them.
    \newline
    A class diagram is the \highlightit{conceptual modeling of the 
    structure of the software application}.
    \newline
    In case of detailed modeling, 
    the class diagram can be \highlightit{translated into programming code}.
    %%% Visibility
    \subsubsection{Visibility}
    \vspace{-0.5cm}
    \begin{table}[ht]
        \begin{tabular}{l l}
            $+$ \quad \quad & public\\
            $-$ & private\\
            \texttt{\#} & protected\\
            \textasciitilde & package\\
        \end{tabular}
    \end{table}
    This notation is placed before methods and attributes names, 
    to indicate their visibility. 
    
    %%% Scope
    \subsubsection{Scope}
    UML defines \highlightit{instance and class scopes for members}.
    \newline
    The latter is represented by \underline{underlined names}.
    \begin{itemize}
        \item \highlightbf{instance members} are scoped to a specific instance
        \item \highlightbf{class members} are the one commonly recognized as 
        \highlighttt{static}
    \end{itemize}
    %%% Arrows
    \subsubsection{Arrows}
    The relationship defined in UML are represented by the following 
    logical connectors
    \begin{figure}[H]
        \centering
        \includegraphics[scale=0.5]{OOPandUML/uml_arrow_cheatsheet}
        \caption{UML logical connectors}
        \label{fig:uml-arrows}
    \end{figure}
    \begin{itemize}
        \item \textbf{Association}\\ 
              a \highlightit{relationship bewteen classes 
              of objects} that allows one object instance to cause another 
              \highlightit{to perform an action on its behalf}.
              \newline
              This relationship is structural and does not represent 
              any behavioural trait.
              \newline
              Code implementation usually requires the requsting object 
              to invoke a method or member function using a 
              \highlightit{reference or a pointer to the controlled 
              object}.
              \newline
              Association can be \highlightit{bidirectional, 
              unidirectional, aggregation and reflexive}.
              \newline
              Objects related via the association are considered to act 
              in a role with respect to the association.
              \newline
              The ends of the association can have all the 
              characteristics of a property: 
              \newline
              \highlightit{multiplicity, name, 
              visibility and a specification} explaining whether the 
              association is ordered and / or unique.
              
        \item \textbf{Dependency}\\ 
              a relationship showing that \highlightit{an element requires 
              other model elements} for their specification or implementation. 
              \newline
              The dashed lines from the dependent/client to 
              the independent/supplier element specifies 
              \highlightit{the direction of a relationship and not of a process}.
              
        \item \textbf{Aggregation}\\
              a variant of \highlightit{has a association relationship 
              and more specific than the latter}. 
              It represents a part-whole or a part-of relationship. 
              As a type of association, an aggregation can be named and 
              have the same adornments an association can.
              \newline
              An aggregation \highlightit{can occur when a class is a 
              collection or a container of other classes}, but 
              \highlightit{the content of the container continues 
              to exist when the container is destroyed}.
              \newline
              In UML notation, the \highlightit{hollow diamond shape is 
              drawn on the containing class end} of the connecting line. 
              \newline
              An example of aggregarion is a Pond which has zero or more 
              Ducks and Duck which has at most one Pond.
              \newline
              \highlightit{A duck can exist separately from a pond}.
              
        \item \textbf{Composition}\\ 
              refers to \highlightit{combining objects within 
              compound objects, ensuring the encapsulation} of each 
              object by \highlightit{using the well-defined interface without 
              exposing the internals}. Differently, data structures do 
              not enforce encapsulation.
              \newline
              Composition can be regarded as a \highlightit{relationship 
              between types}: an object of composite type \highlightit{has} objects 
              of other types
              \newline
              \highlightit{Composition must be distinguished from subtyping}, 
              which is the process of adding detail to a general data type 
              to create a more specific one.
              \newline
              In UML notation, the \highlightit{filled diamond shape is 
              drawn on the containing class end} of the connecting line. An example of 
              composition is a University which has one or more 
              Departments, which belong exactly to one Univerity.
              \newline
              \highlightit{A Department can not exist as a separate entity} 
              detached from a specific University.
    \end{itemize}
    %%% Composition versus Aggregation
    %~ \nparagraph{Composition versus Aggregation}
    %~ \begin{itemize}
        %~ \item Composition relationship
        %~ \item Aggregation relationship
    %~ \end{itemize}
    %%% Class level relationship
    \subsubsection{Class-level relationships}
    \begin{itemize}
        \item \textbf{Generalization/Inheritance}
              \newline
              a relationship between two classes, in which the \highlightit{subclass} is 
              considered to be a specialization from the \highlightit{super type}. 
              \newline
              \highlightit{Any instance of the subtype is also an instance 
              of the superclass}. 
              As an example, \\
              \highlightit{humans are a subclass of simian which is a subclass of mammal}.
              \newline
              Generalization is also known as \highlightbf{inheritance} or a 
              \highlightit{is a} relationship. 
              \newline
              The base class is also known as parent, superclass, 
              base class or base type.
              \newline
              The subtype is also known as child, subclass, 
              derived class, derived type, inheriting class or 
              inheriting type.
        \item \textbf{Realization/Implementation}
              \newline
              a relationship between two model elements, in which
              the \highlightit{client} realizes (implements or executes) the behavior 
              that the \highlightit{supplier} specifies.
    \end{itemize}
    %%% General relationship
    \subsubsection{General relationship}
    \begin{itemize}
        \item \textbf{Dependency}\\
              a weaker form of bond that indicates that one class depends 
              on another because it uses it at some point in time. 
              Sometimes this relationship might be implemented as member 
              function arguments.
        \item \textbf{Multiplicity}\\
              an optional notation which can be 
              useful to add details to the association relationship. At 
              each end of the line, one of the following notation can be 
              used to indicate the range of number of objexts that 
              partecipate in the association from the perspective of the 
              other end.
             \begin{table}[ht]
                 \begin{center}
                     \begin{tabular}{l l}
                         $0$ & No instances\\
                         $0..1$ \quad & No instances, or one instance\\
                         $1$ & Exactly one instance\\
                         $1..1$ & Exactly one instance\\
                         $0..*$ & Zero or more instances\\
                         $*$ & Zero or more instances\\
                         $1..*$ & One or more instances\\
                     \end{tabular}
                 \end{center}
             \end{table}
    \end{itemize}

    \subsubsection{Class and Sequence diagram Cheatsheet}
    Following \highlightit{a summary of the information for class and sequance UML 
    diagrams}
    \begin{figure}[H]
        \centering
        \includegraphics[scale=0.7]{OOPandUML/uml_class_sequence_diag_cheatsheet}
        \caption{Class and Sequence diagrams cheatsheet}
        \label{fig:class-seq-diag-cheatsheet}
    \end{figure}

    \subsection{Sequence Diagram}
    
%%% OOP.
\section{Object Oriented Programming}
    \subsection{Polymorphism}
    In C++ there are different kinds of polymorphism, summarized it 
    the figure~\ref{fig:cpp-poly}
    \begin{figure}[H]
        \centering
        \includegraphics[scale=0.5]{OOPandUML/cpp_polymorphism}
        \caption{Polimorphism in C++}
        \label{fig:cpp-poly}
    \end{figure}
    \subsubsection*{Compile time polymorphism}
    \highlightit{Overloaded functions} are invoked by matching the type 
    and number of arguments at compile time, because the compiler already has all the 
    information needed. This is know by the name of 
    \highlightit{static binding} or \highlightit{early binding}.
    \subsubsection*{Run time polymorphism}

\subsection{Interface vs Abstract class}
\begin{itemize}
    \item Interface can contain only body-less abstract 
    \item Abstract Class
\end{itemize}
